\documentclass[12pt,a4paper]{article}
\usepackage[utf8]{inputenc}
\usepackage[spanish]{babel}
\usepackage{graphicx}
\usepackage{booktabs}
\usepackage{tabularx}
\usepackage[hidelinks]{hyperref}
\usepackage{geometry}
\geometry{top=2.5cm, bottom=2.5cm, left=3cm, right=3cm}

\title{\textbf{Escenario 6: Dise\~no de M\'etricas}}
\author{
\textbf{Asignatura:} Aseguramiento de la Calidad del Software --- NRC 30733 \\
\textbf{Docente:} Ing. Diego Gamboa \\
\textbf{Grupo 7 --- Calidad en Uso} \\
\textbf{Integrantes:} Kevin Amagua\~na, Jairo Bonilla, Darwin Panchez, Eduardo Mortensen
}
\date{\today}

\begin{document}

\maketitle

\section{Descripci\'on General}
Este escenario dise\~na formalmente, siguiendo la estructura de ISO/IEC 25022, las m\'etricas de Calidad en Uso para las tareas priorizadas de MediSalud, estableciendo: f\'ormula, variables, unidad, rango deseado e interpretaci\'on.

\section{Cat\'alogo de M\'etricas ISO/IEC 25022}
A continuaci\'on se presentan las fichas est\'andar para 5 m\'etricas (una por caracter\'istica).

\subsection{M\'etrica 1: Completitud de registro de HCE}
\begin{itemize}
    \item \textbf{Caracter\'istica:} Efectividad
    \item \textbf{Prop\'osito:} Medir el grado en que los m\'edicos logran registrar la nota cl\'inica sin fallos.
    \item \textbf{F\'ormula:} $X = \frac{\text{Notas completadas correctamente}}{\text{Notas intentadas}}$
    \item \textbf{Variables:} $A$ = Notas completadas, $B$ = Notas intentadas
    \item \textbf{Unidad:} Proporci\'on (0-1)
    \item \textbf{Rango deseado:} $X \ge 0.95$
    \item \textbf{Fuente de datos:} Logs de aplicaci\'on HCE
    \item \textbf{Frecuencia de medici\'on:} Semanal
    \item \textbf{Responsable:} Gerencia de Calidad
\end{itemize}

\subsection{M\'etrica 2: Tiempo promedio de registro de HCE}
\begin{itemize}
    \item \textbf{Caracter\'istica:} Eficiencia
    \item \textbf{Prop\'osito:} Evaluar los recursos (tiempo) utilizados por los m\'edicos para completar el registro.
    \item \textbf{F\'ormula:} $X = \frac{\sum t_i}{n}$
    \item \textbf{Variables:} $t_i$ = tiempo invertido por la nota $i$, $n$ = total de notas
    \item \textbf{Unidad:} Segundos
    \item \textbf{Rango deseado:} $X \le 8.0$ segundos (seg\'un RNF-01)
    \item \textbf{Fuente de datos:} Logs con marcas de tiempo (timestamps)
    \item \textbf{Frecuencia de medici\'on:} Continua / Semanal
    \item \textbf{Responsable:} TI / Equipo DevOps
\end{itemize}

\subsection{M\'etrica 3: \'Indice de Satisfacci\'on (CSAT)}
\begin{itemize}
    \item \textbf{Caracter\'istica:} Satisfacci\'on
    \item \textbf{Prop\'osito:} Conocer la percepci\'on de usabilidad del sistema por parte del personal de salud.
    \item \textbf{F\'ormula:} $X = \frac{\sum \text{Puntajes CSAT}}{N \times 5}$
    \item \textbf{Variables:} Puntajes CSAT (escala 1 a 5), $N$ = n\'umero de encuestados
    \item \textbf{Unidad:} Proporci\'on (0-1)
    \item \textbf{Rango deseado:} $X \ge 0.80$
    \item \textbf{Fuente de datos:} Encuestas de satisfacci\'on
    \item \textbf{Frecuencia de medici\'on:} Trimestral
    \item \textbf{Responsable:} Gerencia de Calidad
\end{itemize}

\subsection{M\'etrica 4: Tasa de errores de facturaci\'on}
\begin{itemize}
    \item \textbf{Caracter\'istica:} Libertad de Riesgo
    \item \textbf{Prop\'osito:} Medir el riesgo financiero mitigado en el cobro a seguros.
    \item \textbf{F\'ormula:} $X = \frac{\text{Facturas err\'oneas}}{\text{Facturas emitidas}}$
    \item \textbf{Variables:} Facturas con error reportado, Total de transacciones
    \item \textbf{Unidad:} Proporci\'on
    \item \textbf{Rango deseado:} $X \le 0.01$ (seg\'un RNF-03)
    \item \textbf{Fuente de datos:} Base de datos transaccional e incidentes
    \item \textbf{Frecuencia de medici\'on:} Mensual
    \item \textbf{Responsable:} Departamento de Facturaci\'on
\end{itemize}

\subsection{M\'etrica 5: Consistencia de eficiencia entre sedes}
\begin{itemize}
    \item \textbf{Caracter\'istica:} Cobertura de Contexto
    \item \textbf{Prop\'osito:} Verificar si el sistema rinde igual en todas las sedes (Quito, Guayaquil, etc.).
    \item \textbf{F\'ormula:} $X = 1 - \frac{|\text{Tiempo sede A} - \text{Tiempo sede B}|}{\text{Tiempo sede A}}$
    \item \textbf{Variables:} Tiempo promedio de tarea desagregado por sede
    \item \textbf{Unidad:} Proporci\'on
    \item \textbf{Rango deseado:} $X \ge 0.90$
    \item \textbf{Fuente de datos:} Metadatos de sesi\'on (sede) y logs de tiempo
    \item \textbf{Frecuencia de medici\'on:} Semanal
    \item \textbf{Responsable:} Gerencia de Tecnolog\'ia (TI)
\end{itemize}

\section{Conclusiones}
Se ha traducido conceptos normativos abstractos en m\'etricas formales, reproducibles y accionables, sentando las bases t\'ecnicas para la automatizaci\'on del programa de medici\'on en MediSalud HIS.

\end{document}
